% Diese Seite verwendet eine eigene Seitengeometrie und verzichtet auf die Fußnote.
\newgeometry{left=2cm, right=2cm, top=1cm, bottom=1cm}
\thispagestyle{empty}

% Überschrift der Seite, zentriert
\begin{center}
    \section*{\underline{Symbolverzeichnis}}
\end{center}

\begin{raggedright}
\begin{table}[H]
    \begin{tabular}{l c c l}
        \textbf{Symbol} & \textbf{Einheit} & \textbf{(SI-Einheit)} & \textbf{Bedeutung} \\[2.0ex]
        $A$         & $\left[\text{m}^2\right]$ & & \begin{tabular}[c]{@{}l@{}}Fläche (meistens Querschnitt einer elektrischen Leitung oder\\ Plattenfläche eines Kondensators)\end{tabular} \\[2.5ex]
        $d$         & $\left[\text{m}\right]$ & & Durchmesser (einer Spule oder eines elektrischen Leiters) \\[2.0ex]
        $E$         & $\left[\text{J}\right]$ & $\left[\tfrac{\text{kg}\cdot\text{m}^2}{\text{s}^2}\right]$ & Elektrische Energie \\[2.0ex]
        $\vec{F}$   & $\left[\tfrac{\text{V}}{\text{m}}\right]$ & & Elektrische Feldstärke (oftmals auch $\vec{E}$) \\[2.0ex]
        $G$         & $\left[\tfrac{1}{\mathrm{\Omega}}\right]$ & & Elektrischer Leitwert (Kehrwert zum el. Widerstand)\\[2.0ex]
        $I$, $i(t)$ & $\left[\text{A}\right]$ & & Stromstärke (kleines $i$: zeitabhängig, bspw. bei Wechselstrom) \\[2.0ex]
        $j$         & $\left[-\right]$ & & Komplexe Zahl $j = \sqrt{-1}$ (da $i$ bereits verwendet wird) \\[2.0ex]
        $k$         & $\left[-\right]$ & & Anzahl der Knoten in einem elektrischen Netzwerk \\[2.0ex]
        $m$         & $\left[-\right]$ & & Anzahl der Maschen in einem elektrischen Netzwerk \\[2.0ex]
        $n$         & $\left[-\right]$ & & (Nur als Zähler-Variable genutzt) \\[2.0ex]
        $P$         & $\left[\text{W}\right]$ & $\left[\tfrac{\text{kg}\cdot\text{m}^2}{\text{s}^3}\right]$ & Elektrische Leistung \\[2.0ex]
        $Q$         & $\left[\text{C}\right]$ & $\left[\text{A}\cdot\text{s}\right]$ & Elektrische Ladung \\[2.0ex]
        $r$         & $\left[\text{m}\right]$ & & Radius (einer Spule oder eines elektrischen Leiters) \\[2.0ex]
        $R$         & $\left[\mathrm{\Omega}\right]$ & & Elektrischer Widerstand \\[2.0ex]
        $\rho$      & $\left[\tfrac{\mathrm{\Omega}}{\text{m}}\right]$ & & Spezifischer Widerstand (pro Länge) \\[2.0ex]
        $S$         & $\left[\tfrac{\text{A}}{\text{m}^2}\right]$ & & Stromdichte \\[2.0ex]
        $\sigma$    & $\left[\tfrac{\text{m}}{\mathrm{\Omega}}\right]$ & & Spezifische Leitfähigkeit (Kehrwert zum spez. Widerstand) \\[2.0ex]
        $t$         & $\left[\text{s}\right]$ & & Zeit \\[2.0ex]
        $\vartheta$ & $\left[^\circ\text{C}\right]$ & $\left[\text{K}\right]$ & Temperatur \\[2.0ex]
        $U$, $u(t)$ & $\left[\text{V}\right]$ & & \begin{tabular}[c]{@{}l@{}}Elektrische Spannung (kleines $u$: zeitabhängig, bspw. bei\\ Wechselspannung)\end{tabular} \\[2.5ex]
    \end{tabular}
\end{table}

\end{raggedright}

\restoregeometry